\chapter{Generalità dell'anatomia}

% ---------------------------------------------------------------------------
\section{L'apparato scheletrico}

\textbf{Apparato scheletrico}: funzioni di protezione, sostegno, contributo al movimento, riserva
di minerali ed emopoiesi. Si distingue in:
\begin{itemize}
  \item \textbf{scheletro assile}: cranio, vertebre, coste, sterno, sacro, cartilagini e legamenti
    $\rightarrow$ protegge encefalo, midollo spinale, organi di senso e tessuti molli del torace;
    sostiene il peso corporeo sopra gli arti inferiori;
  \item \textbf{scheletro appendicolare}: arti, ossa di sostegno e legamenti $\rightarrow$ sostegno
    interno e posizione agli arti; sostiene e sposta lo scheletro assile.
\end{itemize}
Lo scheletro = $\sim$20\% del peso corporeo; nell'adulto $\sim$200 ossa. Il \textbf{midollo osseo}
è il sito primario di produzione di globuli rossi e bianchi.

% ---------------------------------------------------------------------------
\section{Struttura dell'osso}

\subsection{Composizione e tipi di tessuto osseo}
Osso = \textbf{matrice extracellulare} dura, compatta e mineralizzata + componente cellulare ricca
di \textbf{osteociti}. Due tipi di tessuto osseo:
\begin{itemize}
  \item \textbf{spugnoso}: rete di lamine e trabecole; riempie le estremità delle ossa lunghe e
    forma lo strato medio delle ossa piatte (es.\ sterno);
  \item \textbf{compatto}: denso e solido, densamente calcificato; forma la superficie esterna di
    tutte le ossa.
\end{itemize}

\subsection{Classificazione in base alla forma}
\begin{itemize}
  \item \textbf{lunghe}: si sviluppano in lunghezza (radio, ulna, omero, tibia, perone);
  \item \textbf{corte}: forma cuboidale (carpo, tarso);
  \item \textbf{piatte}: lunghezza e larghezza $>$ spessore (scapole, sterno, ossa del cranio);
  \item \textbf{irregolari}: non classificabili altrove (vertebre);
  \item \textbf{sesamoidi}: incluse nei tendini (patella, pisiforme).
\end{itemize}

\subsection{Il midollo osseo}
\textbf{Midollo osseo}: tessuto molle nei canali delle ossa lunghe e nella fascia centrale delle
ossa piatte.
\begin{enumerate}
  \item \textbf{giallo}: prevalentemente stroma adiposo (adipociti);
  \item \textbf{rosso}: parenchima emopoietico (cellule del sangue mature e immature, staminali).
\end{enumerate}

\subsection{Vascolarizzazione}
Osso riccamente vascolarizzato e innervato. Due fonti:
\begin{enumerate}
  \item \textbf{arterie nutritizie};
  \item \textbf{arterie periostee}.
\end{enumerate}
La vascolarizzazione diafisaria e metaepifisaria è fittamente anastomizzata.

\subsection{Superficie e periostio}
\textbf{Periostio}: spessa capsula connettivale (tessuto connettivo denso a fibre intrecciate) che
riveste le superfici esterne $\rightarrow$ protegge l'osso e ne supporta l'azione trofica tramite i
vasi di cui è ricco.

Principali caratteristiche di superficie delle ossa:
\begin{table}[htbp]
  \centering \small \renewcommand{\arraystretch}{1.2}
  \begin{tabularx}{\textwidth}{|l|l|X|}
    \hline
    \textbf{Categoria} & \textbf{Termine} & \textbf{Definizione} \\
    \hline
    \multirow{2}{*}{Elevazioni e proiezioni}
      & Processo & proiezione o rilievo \\
      & Ramo & parte che forma un angolo con il resto dell'osso \\
    \hline
    \multirow{6}{*}{Inserzione di tendini/legamenti}
      & Trocantere & voluminosa proiezione rugosa \\
      & Tuberosità & proiezione rugosa più piccola \\
      & Tubercolo & piccola proiezione rotondeggiante \\
      & Cresta & linea rilevata \\
      & Linea & linea sottile \\
      & Spina & processo appuntito \\
    \hline
    \multirow{5}{*}{Articolazione con ossa vicine}
      & Testa & estremità articolare espansa dell'epifisi, separata dalla diafisi da un collo \\
      & Collo & connessione più stretta tra epifisi e diafisi \\
      & Condilo & processo articolare liscio e tondeggiante \\
      & Troclea & processo articolare liscio e scanalato (a puleggia) \\
      & Faccetta & piccola superficie articolare piatta \\
    \hline
    \multirow{2}{*}{Depressioni}
      & Fossa & depressione poco profonda \\
      & Solco & piccola fossa \\
    \hline
    \multirow{4}{*}{Aperture}
      & Forame & foro per vasi e/o nervi \\
      & Fessura & fenditura allungata \\
      & Meato/canale & ingresso a un canale che attraversa l'osso \\
      & Seno/antro & camera interna, in genere contenente aria \\
    \hline
  \end{tabularx}
\end{table}

% ---------------------------------------------------------------------------
\section{Le articolazioni}

\textbf{Articolazioni}: rendono le ossa solidali e ne consentono il movimento reciproco. In base
alla mobilità: \textbf{mobili} (spalla), \textbf{semimobili} (tra vertebre), \textbf{fisse}
(ossa del cranio).

\subsection{Classificazione}
\begin{table}[htbp]
  \centering \small \renewcommand{\arraystretch}{1.2}
  \begin{tabularx}{\textwidth}{|l|l|X|X|}
    \hline
    \textbf{Funzionale} & \textbf{Strutturale} & \textbf{Descrizione} & \textbf{Esempio} \\
    \hline
    \multirow{4}{*}{\shortstack[l]{\textbf{Sinartrosi}\\(nessun mov.)}}
      & Fibrose --- suture & incastro fibroso ad ampie superfici & tra le ossa del cranio \\
      & Fibrose --- gonfosi & inserzioni fibrose nell'alveolo & legamenti periodontali \\
      & Cartilaginee --- sincondrosi & disco cartilagineo interposto & disco epifisario \\
      & Ossee --- sinostosi & ossificazione di un'altra sinartrosi & parti del cranio, linee epifisarie \\
    \hline
    \multirow{2}{*}{\shortstack[l]{\textbf{Amfiartrosi}\\(piccoli mov.)}}
      & Fibrose --- sindesmosi & connessione legamentosa & tra tibia e fibula \\
      & Cartilaginee --- sinfisi & disco fibrocartilagineo & ossa dell'anca, vertebre adiacenti \\
    \hline
    \multirow{3}{*}{\shortstack[l]{\textbf{Diartrosi}\\(ampi mov.)}}
      & Sinoviali --- monoassiale & movimento su un piano & gomito, caviglia \\
      & Sinoviali --- biassiale & movimenti su due piani & coste, polso \\
      & Sinoviali --- triassiale & movimenti su tre piani & spalla, anca \\
    \hline
  \end{tabularx}
\end{table}

\subsection{Le diartrosi sinoviali}
Rivestono le articolazioni mobili. Capsula articolare:
\begin{itemize}
  \item \textbf{membrana fibrosa}: fibre collagene continue col periostio, rinforzata da legamenti
    capsulari;
  \item \textbf{membrana sinoviale}: fissata ai margini della cartilagine articolare; secerne il
    \textbf{liquido sinoviale} (filtrato del plasma) che vascolarizza, nutre e lubrifica la cavità;
    la viscosità è data dall'acido ialuronico.
\end{itemize}

\subsection{Tipi di articolazioni sinoviali}
\begin{itemize}
  \item \textbf{artrodia}: superfici pianeggianti, leggeri scivolamenti (tra processi articolari
    delle vertebre);
  \item \textbf{enartrosi}: capi ``sferici'' (uno concavo, uno convesso), movimenti su tutti i
    piani + rotazione $\rightarrow$ 3 assi, 6 movimenti (coxo-femorale, gleno-omerale);
  \item \textbf{condiloartrosi}: superfici ellissoidali (cavità glenoidea + condilo), movimenti su
    due piani $\rightarrow$ 2 assi, 4 movimenti (polso prossimale). La \textbf{temporo-mandibolare}
    è una diartrosi doppia (disco completo interposto: superiore disco-fossa, inferiore
    disco-condilo);
  \item \textbf{a sella}: superfici a sella di cavallo (concava in un senso, convessa nell'altro),
    2 piani + rotazione assiale (trapezio--I metacarpale, femoro-rotulea);
  \item \textbf{ginglimo angolare (troclea)}: superficie a puleggia, asse $\perp$ alla diafisi,
    movimenti su un piano $\rightarrow$ 1 asse, 2 movimenti (omero-ulnare, femoro-tibiale);
  \item \textbf{ginglimo laterale (trocoide)}: capi cilindrici (uno cavo, uno pieno), asse
    parallelo, movimento rotatorio $\rightarrow$ 1 asse, 2 movimenti (radio-ulnare prossimale,
    atlanto-assiale mediana).
\end{itemize}

\subsection{I movimenti articolari}
\begin{itemize}
  \item \textbf{lineare}: scivolamento senza variazione angolare;
  \item \textbf{angolare}: variazione dell'angolo tra due segmenti;
  \item \textbf{circumduzione}: l'estremità distale descrive un cerchio;
  \item \textbf{rotazione}: attorno all'asse longitudinale.
\end{itemize}
Termini per gli arti: flessione/estensione, abduzione/adduzione, supinazione/pronazione,
eversione/inversione, flessione dorsale/plantare, retrazione/protrusione, opposizione,
depressione/elevazione.

% ---------------------------------------------------------------------------
\section{Orientamento e topografia corporea}

\subsection{Posizione anatomica}
\textbf{Posizione anatomica} (riferimento convenzionale): stazione eretta, viso in avanti, arti
paralleli all'asse del corpo, palmi in avanti (supinazione), piedi leggermente divaricati. Termini
di posizione e piani si definiscono sempre rispetto a questa posizione.

\subsection{Termini di posizione}
\begin{itemize}
  \item \textbf{craniale} (cefalico): vicino al polo cefalico;
  \item \textbf{caudale} (podalico): vicino al polo caudale;
  \item \textbf{mediale}: vicino al piano di simmetria;
  \item \textbf{laterale}: lontano dal piano di simmetria;
  \item \textbf{ventrale}: anteriore;
  \item \textbf{dorsale}: posteriore;
  \item \textbf{prossimale}: vicino all'inserzione (arti);
  \item \textbf{distale}: vicino all'estremità libera (arti).
\end{itemize}

\subsection{Piani e assi di simmetria}
\begin{itemize}
  \item \textbf{piani sagittali}: quello che divide in due parti speculari (\textit{antimeri}) è il
    \textbf{sagittale mediano}; i paralleli sono parasagittali;
  \item \textbf{piani frontali}: andamento cranio-caudale, $\perp$ ai sagittali e ai trasversali;
    separano ventrale da dorsale (nessun piano di riferimento fisso);
  \item \textbf{piani trasversali}: andamento latero-laterale, paralleli alla superficie di
    appoggio; dividono craniale da caudale.
\end{itemize}
Tre assi principali: \textbf{craniocaudale}, \textbf{sagittale}, \textbf{laterolaterale}.

\subsection{Superfici del corpo}
\begin{itemize}
  \item \textbf{corpo}: anteriore (ventrale), posteriore (dorsale), due laterali (dx e sx);
  \item \textbf{arti}: anteriore (ventrale), posteriore (dorsale), laterale, mediale.
\end{itemize}

\subsection{Punti e linee di repere}
\textbf{Punti e linee di repere} (fr.\ \textit{repère} = trovare): strutture scheletriche
apprezzabili dall'esterno $\rightarrow$ localizzano punti e aree rispetto a riferimenti fissi.
Principali repere cervicali e toracici superiori:
\begin{itemize}
  \item processo trasverso dell'atlante;
  \item linea nucale posteriore superiore;
  \item processo trasverso della II vertebra cervicale;
  \item processo spinoso della VII vertebra cervicale;
  \item I costa.
\end{itemize}

% ---------------------------------------------------------------------------
\section{Le cavità del corpo}

\textbf{Cavità del corpo}: contengono e proteggono gli organi interni, permettono i cambiamenti di
dimensione e forma dei visceri.
\begin{itemize}
  \item \textbf{cavità dorsale}:
    \begin{itemize}
      \item \textit{cranica}: encefalo;
      \item \textit{spinale}: midollo spinale;
    \end{itemize}
  \item \textbf{cavità ventrale}:
    \begin{itemize}
      \item \textit{toracica}: mediastino (trachea, bronchi, esofago, timo, cavità pericardica col
        cuore) e cavità pleuriche (polmoni). Il \textbf{diaframma} la separa dall'addominopelvica;
      \item \textit{addominopelvica}:
        \begin{itemize}
          \item \textbf{addominale}: fegato, cistifellea, stomaco, intestino tenue e crasso, milza,
            reni, ureteri;
          \item \textbf{pelvica}: vescica, organi genitali, parte dell'intestino crasso.
        \end{itemize}
    \end{itemize}
\end{itemize}
La cavità addominopelvica si suddivide in \textbf{4 quadranti} (sup./inf.\ dx e sx) oppure in
\textbf{9 regioni}: epigastrica, ipocondriaca dx/sx, ombelicale, lombare dx/sx, ipogastrica,
iliaca dx/sx.

% ---------------------------------------------------------------------------
\section{L'apparato muscolare}

\textbf{Apparato locomotore}: integra 3 sistemi $\rightarrow$ \textbf{scheletrico} (ossa),
\textbf{articolare} (articolazioni), \textbf{muscolare} (muscolatura scheletrica).

\textbf{Muscolo}: tessuto di fibre muscolari con capacità contrattile. Funzioni: movimento dello
scheletro, postura, sostegno a parti molli, protezione, regolazione degli orifizi.

\subsection{Capacità del muscolo}
\begin{itemize}
  \item \textbf{contrattilità}: accorciarsi o allungarsi;
  \item \textbf{eccitabilità}: rispondere a stimoli nervosi e/o ormonali;
  \item \textbf{estensibilità}: essere stirati fino alla lunghezza a riposo dopo la contrazione;
  \item \textbf{elasticità}: tornare alla lunghezza iniziale dopo lo stiramento.
\end{itemize}

\subsection{Tessuto muscolare}
\begin{itemize}
  \item \textbf{striata} (``rossa'', volontaria): regolata dalla volontà; cellule lunghe,
    cilindriche, striate, multinucleate; contrazione attivata da motoneuroni;
  \item \textbf{liscia} (``bianca'', involontaria): attività contrattile autonoma $\rightarrow$
    contrazione prolungata.
\end{itemize}
\textbf{Muscolo cardiaco} (miocardio): istologia simile alla striata volontaria ma funzionamento
involontario, con contrazione ritmica spontanea.

\subsection{Classificazione dei muscoli}
\rubrica{In base al numero di capi (origini)}
\begin{itemize}
  \item \textbf{monocipiti} (1), \textbf{bicipiti} (2), \textbf{tricipiti} (3),
    \textbf{quadricipiti} (4).
\end{itemize}
I \textbf{muscoli scheletrici} hanno origine e inserzione nelle ossa; i \textbf{pellicciai} hanno
almeno un attacco nel derma $\rightarrow$ la contrazione muove la cute.

\rubrica{In base alla forma}
\begin{itemize}
  \item \textbf{lunghi}: lunghezza $>$ larghezza e spessore;
  \item \textbf{larghi}: spessore $\ll$ lunghezza e larghezza;
  \item \textbf{brevi}: lunghezza, larghezza e spessore simili;
  \item \textbf{anulari}: circondano gli orifizi naturali;
  \item \textbf{orbicolari}: si comportano come gli altri scheletrici;
  \item \textbf{sfinteri}: tono accentuato, in continua contrazione.
\end{itemize}

\rubrica{In base all'orientamento delle fibre}
\begin{itemize}
  \item \textbf{unipennati}: fibre su un lato di un tendine centrale (vasto mediale/laterale);
  \item \textbf{semipennati}: fibre su due linee lineari contrapposte (flessore lungo del pollice);
  \item \textbf{bipennati}: fibre convergenti sulle due facce di un tendine centrale (retto del
    femore, gastrocnemio);
  \item \textbf{multi-/pluripennati}: più fasci tendinei con origine comune (deltoide).
\end{itemize}

\rubrica{In base al numero di inserzioni}
\begin{itemize}
  \item \textbf{monocaudati} (1), \textbf{bicaudati} (2), \textbf{tricaudati} (3),
    \textbf{pluricaudati} (più).
\end{itemize}

\rubrica{In base alle iscrizioni tendinee}
\begin{itemize}
  \item \textbf{monogastrici}: nessun tendine intermedio (sopraioidei);
  \item \textbf{digastrici}: un tendine intermedio;
  \item \textbf{poligastrici}: più tendini intermedi (retto dell'addome).
\end{itemize}

\rubrica{In base alla funzione}
Flessori/estensori, adduttori/abduttori, pronatori/supinatori, rotatori interni/esterni. Per ruolo
nel movimento: \textbf{motori primari}, \textbf{sinergici}, \textbf{antagonisti},
\textbf{stabilizzatori}, \textbf{neutralizzatori}.

\subsection{L'angolo di pennazione}
\textbf{Angolo di pennazione}: inclinazione delle fibre rispetto all'asse di trazione del tendine.
\begin{itemize}
  \item \textbf{fasci paralleli} (fusiformi, es.\ bicipite): trasmettono al tendine tutta la
    capacità contrattile;
  \item \textbf{fibre pennate}: ne trasmettono solo una parte $\rightarrow$ a 30° circa l'87\%
    della tensione ($\cos 30^\circ = 0{,}866$).
\end{itemize}
I pennati hanno più fibre compatte dei fusiformi e la \textbf{forza aumenta con l'angolo di
pennazione}. Tre fattori determinano la forza:
\begin{itemize}
  \item \textbf{forma}: grado di contrazione e forza sviluppabile;
  \item \textbf{ipertrofia}: 1 cm\textsuperscript{2} di sezione trasversa $\rightarrow$ $\sim$2--3
    kg di forza;
  \item \textbf{orientamento} delle fibre: ruolo chiave tramite le leve articolari.
\end{itemize}
\textbf{Implicazioni per i calciatori}: un volume eccessivo di quadricipiti e femorali, cresciuto
rapidamente senza adattamento funzionale, può causare gravi infortuni; al limite del volume, i
pennati aumentano l'angolo di pennazione limitando l'espressione di forza (come l'allenamento coi
pesi, che aumenta l'angolo di pennazione).

\subsection{La contrazione muscolare}
\begin{quote}
\textit{«Il muscolo esercita una forza dipendente dal numero delle sue fibre, in rapporto con il
lavoro che può svolgere. Il lavoro muscolare è il prodotto del peso che il muscolo può sollevare
per lo spostamento che può eseguire.»}
\end{quote}
Tre tipi di contrazione:
\begin{itemize}
  \item \textbf{concentrica}: il muscolo si accorcia generando forza (sollevamento del peso);
  \item \textbf{eccentrica}: il muscolo si allunga controllando la forza (abbassamento del peso);
  \item \textbf{isometrica}: si contrae senza accorciarsi, nessuno spostamento articolare.
\end{itemize}

\subsection{Le leve articolari}
\textbf{Leva}: struttura rigida che fa perno su un punto fisso (\textbf{fulcro}). Nelle leve
anatomiche: fulcro = \textbf{articolazione}; \textbf{braccio della forza} = tra fulcro e forza
muscolare; \textbf{braccio del carico} = tra fulcro e resistenza.
\begin{enumerate}
  \item \textbf{primo tipo}: fulcro tra forza e resistenza (testa sul rachide cervicale: muscoli
    posteriori del collo = forza, peso del cranio = resistenza, atlanto-occipitale = fulcro);
  \item \textbf{secondo tipo}: fulcro e forza dallo stesso lato, forza più distante della
    resistenza $\rightarrow$ sempre \textbf{vantaggiosa} (in punta di piedi: fulcro
    metatarso-falangea, resistenza peso corporeo, forza tricipite surale);
  \item \textbf{terzo tipo}: forza più vicina al fulcro della resistenza $\rightarrow$ la più
    comune (flessione dell'avambraccio: fulcro gomito, forza bicipite, resistenza peso
    dell'avambraccio + oggetto).
\end{enumerate}

\subsection{La cellula muscolare}
\textbf{Cellula muscolare} (fibrocellula):
\begin{itemize}
  \item forma \textbf{fusiforme}; si accorcia dopo stimolo;
  \item citoplasma = \textbf{sarcoplasma}; membrana = \textbf{sarcolemma};
  \item contiene \textbf{miofibrille} con le proteine contrattili \textbf{actina} e \textbf{miosina}.
\end{itemize}
